%%%%%%%%%%%%%%%%%%%%%%%%%%%%%%%%%%%%%%%%%
% Important note:
% This template requires the resume.cls file to be in the same directory as the
% .tex file. The resume.cls file provides the resume style used for structuring the
% document.
%
%%%%%%%%%%%%%%%%%%%%%%%%%%%%%%%%%%%%%%%%%

%----------------------------------------------------------------------------------------
%	PACKAGES AND OTHER DOCUMENT CONFIGURATIONS
%----------------------------------------------------------------------------------------

\documentclass{resume}
\usepackage[left=0.75in,top=0.6in,right=0.75in,bottom=0.6in]{geometry}
\usepackage[hidelinks]{hyperref}

\newcommand{\tab}[1]{\hspace{.2667\textwidth}\rlap{#1}}
\newcommand{\itab}[1]{\hspace{0em}\rlap{#1}}

\name{Haoyu Gao}
\address{Melbourne, Victoria, Australia \\ \href{mailto:haoyug1@student.unimelb.edu.au}{haoyug1@student.unimelb.edu.au}}
\address{\href{https://scholar.google.com/citations?user=vSm1V54AAAAJ&hl=en}{Google Scholar} \;\textbullet\; \href{https://orcid.org/0000-0003-3274-8630}{ORCID} \;\textbullet\; \href{https://github.com/Haoyu-Gao}{GitHub} \;\textbullet\; \href{https://haoyu-gao.github.io/}{Homepage}}

\begin{document}

%----------------------------------------------------------------------------------------
%	EDUCATION SECTION
%----------------------------------------------------------------------------------------

\begin{rSection}{Research Interest}
My research spans empirical and automated software engineering, with a focus on how developers create, maintain, and use documentation. I study documentation as a socio-technical artefact, combining repository mining, qualitative methods, and machine learning to analyse both its technical content and the human practices around it. My current work examines how development knowledge is communicated to AI agents. I am particularly interested in agent configuration files as a medium for encoding developer intent, context, and process knowledge to support effective human–AI collaboration.
    
\end{rSection}




\begin{rSection}{Education}
%--copy and paste this region  if you need more--
{\bf The University of Melbourne} \hfill {\em Jan 2023 - Nov 2026 (Expected)} 
\\ Doctor of Philosophy - Engineering and IT \hfill {}
\\ Supervisors: A/Prof. Christoph Treude and Dr. Mansooreh Zahedi\\
Research Topic: \textit{Understanding and Maintaining Documentation in Fast-Evolving Software Ecosystems.}
%--copy and paste this region  if you need more--

{\bf The University of Melbourne} \hfill {\em Feb 2021 - Dec 2022} 
\\ Master of Information Technology \hfill {WAM: 86.7/100} 
% Coursework: Multivariate Statistics, Cluster and Cloud Computing, Software Process Management, Statistical Machine Learning, Natural Language Processing, Computer Vision.

{\bf Fuzhou University} \hfill {\em Sep 2016 - July 2020} 
\\ Bachelor of Science - Mathematics and Applied Mathematics \hfill {WAM: 86.6/100} 
% Coursework: Mathematics Analysis, Advanced Algebra, Real Analysis, Complex Analysis, Function Analysis, Ordinary Differential Equation, Partial Differential Equation, Operation Research, Convex Optimisation, Modern Algebra, Topology, Differential Geometry, Fuzzy Mathematics, Graph Theory.



\end{rSection}
%----------------------------------------------------------------------------------------
%	EXPERIENCE SECTION
%----------------------------------------------------------------------------------------

\begin{rSection}{Publications}

    [TSE'25] \textbf{H. Gao}, C. Treude, and M. Zahedi. ``Adapting Installation Instructions in Rapidly Evolving Software Ecosystems''. In \textit{IEEE Transactions on Software Engineering}, 2025.  \textbf{(core A*)}
    

    [FSE'23] \textbf{H. Gao}, C. Treude, and M. Zahedi. ``Evaluating Transfer Learning for Simplifying GitHub READMEs''. In \textit{ESEC/FSE'23: Proceedings of the Joint European Software Engineering Conference and Symposium on the Foundations of Software Engineering, 2023.} \textbf{(core A*)} 

     [EMSE'26] \textbf{H. Gao}, M. Zahedi, W. Jiang, H. Lin, J. Davis, and C. Treude. ``AI Failures in the Eyes of the Downstream Developer: A First Look at Concerns, Practices, and Challenges''. In \textit{Empirical Software Engineering (Springer Nature)}, 2026. \textbf{(core A)}
    
    
    [ESEM'24] \textbf{H. Gao}, M. Zahedi, C. Treude, S. Rosenstock, and M. Cheong. ``Documenting Ethical Considerations in Open Source AI Models''. In \textit{ESEM'24: 18th International Symposium on Empirical Software Engineering and Measurement}, 2024.  \textbf{(core A)}




    [ACL'25 findings] H. Lin, C. Liu, \textbf{H. Gao}, P. Thongtanunam, and C. Treude. CodeReviewQA: The Code Review Comprehension Assessment for Large Language Models. In \textit{Findings of ACL'25}: The 63rd Annual Meeting of the Association for Computational Linguistics \textbf{(core A*)}
    

    [ICSME'25] P. Banyongrakkul, M. Zahedi, P. Thongtanunam, C. Treude, and \textbf{H. Gao}. From Release to Adoption: Challenges in Reusing Pre-trained AI Models for Downstream Developers. In \textit{ICSME'25: 41st International Conference on Software Maintenance and Evolution} \textbf{(core A)}

    [MSR'26 short] \textbf{H. Gao}, P. Banyongrakkul, H. Guan, M. Zahedi, C. Treude, ``On Autopilot? An Empirical Study of Human-AI Teaming and Review Practices in Open Source''. \textit{Short Paper} in \textit{MSR'26: 23rd International Conference on Mining Software Repositories}, 2026, \textbf{(core A)} 
    


    \smallskip
    \hrule
    \smallskip
    {\bf \textit{Under Review \& Preprints}}
    \smallskip

    \textbf{H. Gao}, H. Lin, C. Treude, G. Gay, and M. Zahedi. Does My README File Need To Be Updated? Exploring LLM-Based README Maintenance. (\textbf{Major revision under TSE)}

    H. Lin, C. Liu, \textbf{H. Gao}, P. Thongtanunam, and C. Treude. Fine-grained Approaches for Confidence Calibration of LLMs in Automated Code Revision. \textbf{(Major revision under TSE)}
    
     P. Banyongrakkul, M. Zahedi, C. Treude, \textbf{H. Gao} and P. Thongtanunam. When AI Models Become Dependencies: Studying the Evolution of Pre-Trained Model Reuse in Downstream Software Systems. (\textbf{Under submission})

    \textbf{H. Gao}, J. Lulla, H. Lin, S. Baltes, C. Treude, M. Zahedi. From Registry to Repository: How AI Agent Skills Are Written, Adapted, and Maintained. (\textbf{Under submission})
     

    L. Salerno, \textbf{H. Gao}, P. Thongtanunam, C. Treude. ChatGPT as an Installation Assistant: How Novices Use LLMs for Software Tool Installation. \textbf{(Under submission)}

\end{rSection}





\begin{rSection}{Experience}

{\bf Visiting Postgraduate Research Student}{, Singapore Management University}\hfill {\em Sep 2025- Feb 2026}

Host Professor: A/Prof. Christoph Treude



{\bf Teaching Assistant}{, The University of Melbourne} \hfill {\em Oct 2023- Present}

\textit{Tutored Courses}:
\begin{itemize}
    \item COMP90041 Programming and Software Development
    \item SWEN90017 Masters Advanced Software Project
\end{itemize}

\textit{Marked Courses}:
\begin{itemize}
    \item COMP90041 Programming and Software Development
    \item SWEN90016 Software Processes and Management
    \item COMP90049 Introduction to Machine Learning
    \item COMP90051 Statistical Machine Learning
\end{itemize}

% \begin{itemize}
%     \item Tutored Course: COMP90041 Programming and Software Development
%     \item Marked Courses: COMP90041 Programming and Software Development, SWEN90016 Software Processes and Management, COMP90049 Introduction to Machine Learning, COMP90051 Statistical Machine Learning
% \end{itemize}
Skills: Java Programming Language, Object-Oriented Programming, Software Development Processes, Agile Development, Machine Learning, Assignment Design, University Teaching.


{\bf Research Assistant}{, The University of Melbourne} \hfill {\em Dec 2023- May 2024}\\
Skills: Mining software repositories, Qualitative analysis, Ethics in Machine Learning

{\bf Machine Learning Internship}{, Haier Group Corporation} \hfill {\em Oct 2021- Feb 2022}\\
Skills: Machine Learning, Data analysis

\end{rSection}
%--------------------------------------------------------------------------------
%    PROJECTS
%-----------------------------------------------------------------------------------------------
% \begin{rSection}{Projects}
% %--copy and paste this region  if you need more--
% {\bf Title of the project}{, institution associated with it} \hfill {\em start date- end date}\\
% description of the project

% \end{rSection}
%--------------------------------------------------------------------------------
%    ACTIVITIES
%-----------------------------------------------------------------------------------------------
\begin{rSection}{Services and Activities}


\textbf{Journal Review}
\begin{itemize}
    \item Empirical Software Engineering Journal.
    \item Automated Software Engineering Journal.
\end{itemize}

\textbf{Conference Review}
\begin{itemize}
    \item PC member for ICSME'26 Demo and Data Showcase Track.
    \item PC member for ICSE'26 Artifact Evaluation Track.
    \item Junior PC member for MSR'25.

\end{itemize}

\textbf{Sub-reviewer}\\
TOSEM, ICSE'25, ESEM'24, APSEC'24.

\textbf{Student Volunteer} \\
Event: ICSE 2023 - International Conference on Software Engineering\\ 
Skills: registration assistance and presentation session support. 

%--copy and paste this region  if you need more--
\end{rSection}




\begin{rSection}{(Co)-Supervised Students}
    Hao Guan - Master of Software Engineering, The University of Melbourne
\end{rSection}


\begin{rSection}{Skills}
{\bf Programming}
\\Python, Java, JavaScript, R.

{\bf Research}
\\Data mining, Interview-based study, Qualitative analysis, Machine learning, HPC cluster.

% {\bf Research Methodology}
% \\Quantitative Research, Qualitative Research.\\
{\bf Languages}
\\ Mandarin (Native), English (Full professional proficiency).\\
\end{rSection}



\begin{rSection}{Awards and Scholarships} 

{\bf FEIT Conference Travel Grant}{, The University of Melbourne
} \hfill{\em Feb 2026}
\\ Funded by the faculty and school for conference traveling with A\$4,000.

{\bf 2022 Dean's Honours List}{, The University of Melbourne
} \hfill{\em July 2023}
\\The Dean’s Honours List recognises outstanding academic performance in the academic year, representing the top five
per cent of students enrolled in an engineering or information technology master by coursework program. 

{\bf Melbourne Research Scholarship}{, The University of Melbourne
} \hfill{\em Jan 2023}
\\This scholarship is rewarded to high-achieving students participating in research activities.

{\bf Melbourne Graduate Scholarship}{, The University of Melbourne
} \hfill{\em Jan 2022}
\\This scholarship is rewarded to coursework students in recognition of their academic achievements.

\end{rSection}

% \begin{rSection}{Conference Attendance}
    
% \end{rSection}









\end{document}-------------------

